\section{The exact period-eight phase}
\label{sec:period8}

Theorem~\ref{thm:flux-chiral} explains a dimension-halving mechanism but does
not solve the resulting block at a general period.  The word $\tau_*$ is the
first sub-eight phase and has enough additional structure to reduce its
$4\times4$ squared block once more.  This section carries that calculation
through to the complete dispersion and then returns to the finite rings.

\subsection{The target fiber and its chiral block}

For $|z|=1$, the period-eight fiber in the residue order $0,1,\ldots,7$ is
\begin{equation}
 H(z)=\begin{pmatrix}
0&1&1&0&0&0&z^{-1}&z^{-1}\\
1&0&1&1&0&0&0&-z^{-1}\\
1&1&0&1&-1&0&0&0\\
0&1&1&0&1&1&0&0\\
0&0&-1&1&0&1&-1&0\\
0&0&0&1&1&0&1&-1\\
z&0&0&0&-1&1&0&1\\
z&-z&0&0&0&-1&1&0
\end{pmatrix}.
\label{eq:period8-fiber}
\end{equation}
The boundary entries occur in conjugate pairs, so $H(z)$ is Hermitian.  Since
$\tau_{i+4}=-\tau_i$, Section~\ref{sec:chiral} applies.  Choose $\xi$ with
$\xi^2=z$.  Here $m=4$, so $\gamma_4(z)=\xi^{-1}$, and the normalized signed
four-site shift is the chiral involution $J_z$.

In bases of its two eigenspaces, $H(z)$ is off diagonal.  Multiplying the two
off-diagonal blocks gives the following squared block:
\begin{equation}
 S(\xi)=
 \begin{pmatrix}
4-h&0&1+\xi^{-1}&2\\
0&4-h&2&1-\xi^{-1}\\
1+\xi&2&4+h&0\\
2&1-\xi&0&4+h
 \end{pmatrix},
 \qquad h=\xi+\xi^{-1}.
 \label{eq:period8-squared-block}
\end{equation}
The squared eigenvalues of $H(z)$ are the eigenvalues of $S(\xi)$, each
occurring in both chiral blocks.  Replacing $\xi$ by $-\xi$ exchanges the two
chiral coordinate spaces; the final characteristic data depend only on $z$.

\subsection{The period-eight polynomial}

Write \eqref{eq:period8-squared-block} in $2\times2$ blocks as
\[
 S(\xi)=
 \begin{pmatrix}(4-h)I_2&U(\xi)\\V(\xi)&(4+h)I_2\end{pmatrix},
\]
where
\[
 U(\xi)=\begin{pmatrix}1+\xi^{-1}&2\\2&1-\xi^{-1}\end{pmatrix},
 \qquad
 V(\xi)=\begin{pmatrix}1+\xi&2\\2&1-\xi\end{pmatrix}.
\]
Their product has the form
\begin{equation}
 U(\xi)V(\xi)=
 \begin{pmatrix}
 6+h&4+2(\xi^{-1}-\xi)\\
 4+2(\xi-\xi^{-1})&6-h
 \end{pmatrix},
 \quad \tr(UV)=12,\quad \det(UV)=10+3c.
 \label{eq:uv-invariants}
\end{equation}
Because the diagonal blocks are scalar, the elementary block determinant
identity gives, without an exceptional Schur-complement value,
\begin{equation}
 \det(yI_4-S(\xi))
 =\det\!\left(\bigl((y-4)^2-h^2\bigr)I_2-U(\xi)V(\xi)\right).
 \label{eq:two-by-two-determinant}
\end{equation}
Using $h^2=c+2$, the determinant in
\eqref{eq:two-by-two-determinant} is $q^2-12q+10+3c$, where
$q=(y-4)^2-h^2$.  Expanding this scalar quadratic yields
\begin{equation}
\begin{split}
 P(y,c)={}&y^4-16y^3+(80-2c)y^2\\
          &+(-128+16c)y+c^2-13c+38.
\end{split}
\label{eq:period8-polynomial}
\end{equation}
Thus
\begin{equation}
 \det(\lambda I_8-H(z))=P(\lambda^2,z+z^{-1}).
 \label{eq:period8-characteristic}
\end{equation}
The Bloch phase enters the characteristic polynomial only through the real
variable $c\in[-2,2]$.

\subsection{Complete dispersion}

Put $y=X+4$.  The quartic becomes even in $X$:
\begin{equation}
 P(X+4,c)=X^4-(16+2c)X^2+c^2+19c+38.
 \label{eq:centered-quartic}
\end{equation}
With $W=X^2$, the two roots are
\begin{equation}
 W_\delta(c)=8+c+\delta\sqrt{26-3c},\qquad \delta\in\{\pm1\}.
 \label{eq:w-branches}
\end{equation}
For $-2\leq c\leq2$, both are positive.  Indeed,
\[
 (8+c)^2-(26-3c)=c^2+19c+38\geq4,
\]
and $8+c>0$.  The four squared branches are therefore
\begin{equation}
 y_{\varepsilon,\delta}(c)=4+\varepsilon\sqrt{W_\delta(c)},
 \qquad \varepsilon,\delta\in\{\pm1\}.
 \label{eq:four-squared-branches}
\end{equation}
Since $0<W_-<W_+<16$, their order is uniform:
\[
 y_{+,+}>y_{+,-}>y_{-,-}>y_{-,+}>0.
\]
Moreover,
\[
 W_+'(c)=1-\frac{3}{2\sqrt{26-3c}}>0,
 \qquad
 W_-'(c)=1+\frac{3}{2\sqrt{26-3c}}>0.
\]
Thus the two upper branches increase with $c$, while the two lower branches
decrease.  In particular the top squared edge is
\begin{equation}
 r(c)=4+\sqrt{8+c+\sqrt{26-3c}},
 \label{eq:top-dispersion}
\end{equation}
strictly increasing on $[-2,2]$.

Set
\[
 a=\sqrt{10+2\sqrt5},\qquad b=\sqrt{10-2\sqrt5}.
\]
At $c=-2$ the four squared roots are
\[
 6+\sqrt2,\quad6-\sqrt2,\quad2+\sqrt2,\quad2-\sqrt2,
\]
whereas at $c=2$ they are $4+a,4+b,4-b,4-a$.  The four squared bands are
therefore
\begin{equation}
\begin{array}{ll}
 [6+\sqrt2,\,4+a],&[6-\sqrt2,\,4+b],\\
 {}[4-b,\,2+\sqrt2],&[4-a,\,2-\sqrt2].
\end{array}
\label{eq:squared-bands}
\end{equation}
The middle separation has length $4-2\sqrt2>0$.  The other two are equal to
$2+\sqrt2-b>0$; the last inequality follows after squaring from
$10-2\sqrt5<6+4\sqrt2$.  Hence the four intervals are disjoint.  Equation
\eqref{eq:period8-characteristic} then recovers
the eight simple Hermitian eigenvalues
\[
 \{\,\pm\sqrt{y_{\varepsilon,\delta}(c)}:
       \varepsilon,\delta\in\{\pm1\}\,\}.
\]
In particular, the spectrum has the exact central gap
\begin{equation}
 \left(-\sqrt{4-a},\sqrt{4-a}\right).
 \label{eq:central-gap}
\end{equation}

The substitution $y=X+4$ reveals more than a convenient solution of the
quartic: all odd powers of $X$ disappear.  This second algebraic symmetry is
special to the period-eight block and goes beyond the dimension halving forced
by chirality alone.

\subsection{Finite phase quantization}

The infinite unit-circle maximum occurs uniquely at $c=2$, or $z=1$, and is
\[
 r(2)=4+a=\etaedge.
\]
For positive holonomy the finite grid $z^L=1$ always contains $z=1$.
Proposition~\ref{prop:finite-bloch} therefore gives both an upper bound and an
attaining fiber:
\begin{equation}
 \rho(A_{8L,+}^{*})^2=\etaedge
 \qquad(L\geq1).
 \label{eq:positive-finite-exact}
\end{equation}
The two edge eigenvalues $\pm\sqrt\etaedge$ are simple in $H(1)$ because the
four squared roots at $c=2$ are distinct.

For negative holonomy, the allowed phases are
\[
 z_j=\exp\!\left(\frac{(2j+1)\pi\mathrm i}{L}\right).
\]
The largest value of $z+z^{-1}$ is $2\cos(\pi/L)$.  Strict monotonicity of
$r$ gives the second exact finite formula
\begin{equation}
 \rho(A_{8L,-}^{*})^2
 =4+\sqrt{8+2\cos(\pi/L)+
                \sqrt{26-6\cos(\pi/L)}}.
 \label{eq:negative-finite-exact}
\end{equation}
It is strictly below $\etaedge$ for every finite $L$ and increases to
$\etaedge$ as $L\to\infty$.  Both holonomy sectors are therefore samples of
one dispersion law; their difference is entirely the finite phase grid.

\subsection{The twisted comparison}
\label{sec:twisted}

The natural twisted signing provides the comparison relevant to
\eqref{eq:fixed-minimum}.  In
Hamilton-gauge coordinates, this is the alternating triangle-flux class with
negative Hamilton holonomy.  Its two-site Fourier block is
\begin{equation}
 B(\theta)=2\begin{pmatrix}
 \cos\theta&\cos2\theta\\
 \cos2\theta&-\cos\theta
 \end{pmatrix},
 \label{eq:twisted-block}
\end{equation}
so its eigenvalues have squared magnitude $4g(\theta)$, where
$g(\theta)=\cos^2\theta+\cos^22\theta$.  The shifted finite grid is
$\theta=(2k+1)\pi/(8L)$.  By
$g(-\theta)=g(\theta)=g(\pi-\theta)$ it is enough to consider
$\theta\in[\pi/(8L),\pi/2]$.  With $u=\cos^2\theta$,
\[
 g(\theta)=4u^2-3u+1.
\]
Thus $g(\theta)\leq1$ on $[\pi/6,\pi/2]$.  On $(0,\pi/6]$,
\[
 g'(\theta)=-\sin(2\theta)(1+4\cos(2\theta))<0.
\]
Finally $g(\pi/(8L))\geq g(\pi/8)=(4+\sqrt2)/4>1$.  Hence the first positive grid
phase is maximal, and
\begin{equation}
 \rho(A_{8L}^{\mathrm{tw}})^2
 =4+2\cos\frac{\pi}{4L}+2\cos\frac{\pi}{2L}.
 \label{eq:twisted-radius}
\end{equation}

\begin{proof}
The expression in \eqref{eq:twisted-radius} is increasing for $L\geq4$, so
comparison with \eqref{eq:positive-finite-exact} reduces to $L=4$.  The
half-angle identities give
\begin{equation}
 2\cos\frac{\pi}{16}=\sqrt{2+\sqrt{2+\sqrt2}},
 \qquad 2\cos\frac{\pi}{8}=\sqrt{2+\sqrt2}.
 \label{eq:rational-separator}
\end{equation}
The three rational bounds
\[
 \sqrt5<\frac{179121}{80000},\qquad
 \sqrt{2+\sqrt2}>\frac{923}{500},
 \qquad
 \sqrt{2+\sqrt{2+\sqrt2}}>\frac{1961}{1000}
\]
follow by squaring positive quantities.  The first gives
$\etaedge<1561/200$, while the other two make the $L=4$ twisted value larger
than $4+3807/1000>1561/200$.  This proves the strict comparison for every
$L\geq4$.  The conclusion for $m(C_{8L}(1,2))$ follows by using
$A_{8L,+}^{*}$ as an admissible signing.
\end{proof}

The separator appears only in this comparison of two exact values.  The
original $\alpha=+1$ finite comparison also admits a supplementary Lean
verification.

The period-eight calculation explains both solvability and finite attainment.
It does not select the period: a shorter chiral word might conceivably have
crossed the same threshold.  Nor does it decide whether $Q_*$ is isolated
among period-eight words.  Both questions concern defect geometry rather than
the solved quartic.
